\section{Logic trên đồ thị}
\begingroup\small
\begin{multicols}{2}
\subsection{Mệnh đề và đồ thị kéo theo}
Trong 2-SAT, mỗi điều kiện là một mệnh đề có tối đa hai literal,
ví dụ $(a\lor b)$. Một literal là biến $x$ hoặc phủ định
$\neg x$. Mệnh đề tương đương hai kéo theo
$\neg a\Rightarrow b$ và $\neg b\Rightarrow a$.
Biểu diễn mỗi literal bằng một đỉnh; phủ định là đỉnh ghép cặp.
Một đường đi $p\leadsto q$ nghĩa là nếu chọn $p$ đúng thì
buộc $q$ đúng. Với một yêu cầu đơn $a$, thêm $(a\lor a)$.
Khi thêm yêu cầu $a\Rightarrow b$, phải thêm cả cạnh
$a\to b$ và cạnh phản đảo $\neg b\to\neg a$;
hàm \texttt{implication(a,b)} bên dưới thực hiện đủ hai cạnh.
Với yêu cầu $a\Leftrightarrow b$, thêm hai chiều kéo theo;
với yêu cầu đúng đúng một trong hai, thêm
$(a\lor b)\land(\neg a\lor\neg b)$.

\subsection{Bài toán 2-SAT}
Tạo đồ thị kéo theo $2n$ đỉnh, tìm các thành phần liên thông
mạnh (SCC). Công thức có nghiệm khi và chỉ khi không có biến
$x$ nào mà $x$ và $\neg x$ cùng SCC: nếu cùng SCC, chọn một
giá trị sẽ kéo theo phủ định của chính nó. Chiều ngược lại
suy ra bằng cách xét DAG các SCC theo thứ tự tô pô và chọn
một trong mỗi cặp SCC phủ định. Độ phức tạp $O(n+m)$ với
$m$ mệnh đề. Ràng buộc có ba literal trở lên không được đưa
thẳng vào mô hình 2-SAT; SAT tổng quát khó hơn về bản chất.

\subsection{Dựng và đọc một nghiệm}
Trong mẫu code, biến đánh số $i\in[0,n)$; dùng đỉnh $2i$
cho $\neg x_i$, $2i+1$ cho $x_i$,
nên phủ định một literal là XOR với $1$. Kosaraju tạo số SCC
tăng dần theo thứ tự tô pô của DAG SCC; đặt
$x_i=[\operatorname{comp}(x_i)>
\operatorname{comp}(\neg x_i)]$.
Quy ước so sánh có thể đảo nếu đổi cách đánh số SCC, vì thế
nên kiểm bằng mệnh đề đơn $(x)$ và $(\neg x)$.
\texttt{solve()} trả rỗng khi vô nghiệm, ngược lại là vector
$0/1$ của các biến.

\paragraph{Mẫu 2-SAT (C++17).}
\cppsnippet{code/graph_logic/two_sat.hpp}
\end{multicols}
\endgroup
